
% =============================================================================
% Project : THE_PHILOSOPY_OF_VIABILITY_LIBRARY_PROJECT
% Book    : VS.B01_ONTOLOGY_OF_VIABILITY
% File    : VS.B01_ONTOLOGY_OF_VIABILITY/glossary.tex
% Location: VS.B01_ONTOLOGY_OF_VIABILITY/
%
% Purpose:
%   Provides the controlled philosophical glossary for the first book
%   of The Philosophy of Viability.
%
% Editorial principles:
%   1. Glossary entries clarify terminology without replacing the
%      philosophical arguments developed in the chapters.
%   2. Established principles are distinguished from provisional
%      terminology and unresolved questions.
%   3. No entry should introduce an ontological commitment that has
%      not been established by the philosophical framework.
%   4. The glossary is a controlled terminological reference and may
%      evolve as the philosophical system develops.
%   5. Definitions must remain consistent with the architecture of
%      the four parts of this book.
%
% Dependencies:
%   Standard LaTeX commands only.
% =============================================================================


% -----------------------------------------------------------------------------
% Glossary Heading
% -----------------------------------------------------------------------------

\chapter*{Glossary of Philosophical Terms}
\addcontentsline{toc}{chapter}{Glossary of Philosophical Terms}
\markboth{Glossary of Philosophical Terms}
         {Glossary of Philosophical Terms}


This glossary provides controlled working definitions of central terms
used throughout \textit{Ontology of Viability}.

The entries are intended to establish terminological consistency across
the book. They do not replace the arguments through which the relevant
concepts are examined, differentiated, or justified.

The glossary distinguishes between established conceptual commitments,
structural terminology, and terms whose final philosophical determination
remains open.


% =============================================================================
% PART I — ESTABLISHED ONTOLOGICAL TERMINOLOGY
% =============================================================================

\section*{I. Established Ontological Terminology}


\subsection*{Being}

The ontological condition or reality of that which is.
Being is not treated as a merely linguistic category or as something
whose existence depends on being known by a subject.


\subsection*{Existence}

The actual being of something.
Within the framework of this philosophy, existence and identity are
ontologically inseparable: to exist is to have identity, and identity
is not an additional property subsequently attributed to an otherwise
indeterminate existent.

This inseparability does not establish a temporal or logical priority
of either term over the other.


\subsection*{Identity}

The ontological determination through which an existent is identifiable
as what it is.

Identity is not understood as an external label or as an attribute
added to an independently existing, initially neutral entity.

The precise ontological constitution and dynamics of identity are
developed in the relevant parts of this book.


\subsection*{Difference}

A determination through which distinction becomes intelligible or
ontologically relevant.

Difference must not automatically be equated with the negation of
identity. The relationship between identity and difference requires
examination within the broader ontological framework.


\subsection*{Determination}

The process or condition through which an existent is specified rather
than remaining merely indeterminate in description.

The precise forms, levels, and ontological implications of determination
are examined throughout the book.


\subsection*{Manifestation}

The manner or condition in which being becomes manifest.

Manifestation is not necessarily equivalent to complete disclosure,
perceptual accessibility, or exhaustive knowledge of what is manifested.


\subsection*{Reality}

That which has ontological standing independently of whether, and to
what extent, it is known by a knowing subject.

Reality is not reducible to the subject's representations, descriptions,
or available knowledge.


\subsection*{Knowing Subject}

An existent engaged in the activity or condition of knowing.

The knowing subject is situated within being and is not treated as an
external observer standing outside reality as a whole.


\subsection*{Knowledge}

A mode of relation between a knowing subject and that which is known.

Knowledge is situated and domain-dependent. Knowledge in one domain
may coexist with lack of knowledge in other domains.


\subsection*{Ignorance}

The absence, limitation, or incompleteness of knowledge concerning
something.

Ignorance is not intrinsically a negative value or a philosophical
defect. It may coexist with genuine knowledge and is part of the
situated condition of the knowing subject.


% =============================================================================
% PART II — VIABILITY AND ITS ONTOLOGICAL CONTEXT
% =============================================================================

\section*{II. Viability and Its Ontological Context}


\subsection*{Viability}

The central concept of the philosophical system of viability.

Viability must not be reduced in advance to biological survival,
practical success, persistence, adaptation, or the achievement of
a predefined objective.

Its full ontological meaning and its relation to being, identity,
determination, and dynamics are matters for the systematic development
of the philosophy.


\subsection*{The Philosophy of Viability}

The philosophical system whose foundational ontological framework
is developed in this book.

It investigates being, identity, manifestation, the realms of being,
ontological light, and the dynamics of being as interconnected
dimensions of a broader philosophical architecture.


\subsection*{Ontological Foundation}

The foundational account of being and its constitutive conditions
within which further philosophical investigations can be situated.

An ontological foundation is not interchangeable with an epistemological
foundation or a methodological starting point.


\subsection*{Ontological Independence}

The recognition that reality possesses ontological standing that is
not exhausted by, or dependent upon, the knowledge available to
a knowing subject.

This does not imply that the knowing subject exists outside being.


\subsection*{Ontological Relation}

A relation considered in terms of its place or significance within
being itself, rather than exclusively as a relation established
by thought, language, or observation.


% =============================================================================
% PART III — STRUCTURAL TERMINOLOGY OF THE REALMS OF BEING
% =============================================================================

\section*{III. Structural Terminology of the Realms of Being}


\subsection*{Realm of Being}

A differentiated domain within the structural account of being.

The realms are not to be interpreted automatically as successive
temporal stages, independent worlds, or equivalent hierarchical levels.


\subsection*{Fundamental Realms}

The four structurally fundamental realms identified in the current
architecture of the philosophy:

\begin{itemize}
    \item First Determination
    \item Structural Determination
    \item Relational Manifestation
    \item Material Manifestation
\end{itemize}

Their precise ontological distinctions and relations are developed
in Part II of this book.


\subsection*{First Determination}

The first of the four fundamental realms in the current structural
architecture.

Its precise ontological constitution and relation to the other
fundamental realms are to be established through the corresponding
chapters.


\subsection*{Structural Determination}

The second fundamental realm in the current structural architecture.

It identifies a distinct domain of determination whose full
ontological meaning must be understood through its systematic
development.


\subsection*{Relational Manifestation}

The third fundamental realm in the current structural architecture.

It identifies a distinct domain of manifestation associated with
relational determination. Its exact scope and ontological function
are developed in Part II.


\subsection*{Material Manifestation}

The fourth fundamental realm in the current structural architecture.

It identifies a distinct domain of manifestation associated with
material determination. Its precise meaning must not be reduced
to an assumed physicalist interpretation without further argument.


\subsection*{Intermediate Realms}

The three intermediate or transitional realms identified in the
current architecture:

\begin{itemize}
    \item Hub of Decree
    \item Hub of Signification
    \item Hub of Realization
\end{itemize}

They are not to be treated as three additional fundamental realms
of the same structural status.

Their exact relations to the four fundamental realms require
systematic examination.


\subsection*{Hub of Decree}

The first of the three intermediate realms in the current architecture.

Its precise ontological role, scope, and relation to the fundamental
realms remain subject to systematic development.


\subsection*{Hub of Signification}

The second of the three intermediate realms in the current architecture.

Its precise ontological role and relation to signification, determination,
and manifestation require further elaboration.


\subsection*{Hub of Realization}

The third of the three intermediate realms in the current architecture.

Its precise ontological role and relation to realization and the
fundamental realms remain subject to systematic development.


\subsection*{Transcendent Ground}

A term associated with the question of a possible transcendent
ontological ground.

Its precise status and relation to the structural architecture of
the realms have not been established as a definitive commitment
in this glossary.


\subsection*{Throne-Domain}

A term whose possible ontological status and relation to the realms
of being remain open.

This glossary does not assign it a definitive position within the
architecture of the fundamental or intermediate realms.


% =============================================================================
% PART IV — LIGHT AND MANIFESTATION
% =============================================================================

\section*{IV. Light and Manifestation}


\subsection*{Light}

A central ontological concept developed in Part III of this book.

Light must not be assumed to be identical with epistemic illumination,
human understanding, perceptual visibility, or the acquisition
of knowledge.

Its ontological meaning requires independent examination.


\subsection*{Ontological Light}

Light considered in its ontological significance.

It must be distinguished from light understood merely as a metaphor
for knowledge, awareness, or intellectual clarity.


\subsection*{Epistemic Light}

Light understood in relation to knowing, understanding, disclosure,
or epistemic accessibility.

Epistemic light and ontological light are distinguishable concepts
and must not be conflated.


\subsection*{Disclosure}

The becoming accessible or manifest of something in a particular
respect.

Disclosure does not necessarily entail exhaustive knowledge or
the complete manifestation of that which is disclosed.


% =============================================================================
% PART V — DYNAMICS OF BEING
% =============================================================================

\section*{V. Dynamics of Being}


\subsection*{Dynamics of Being}

The domain of inquiry concerned with the dynamic dimensions of being.

It is not restricted to observable change, physical movement, or
temporal succession.


\subsection*{Change}

A difference or transformation in the condition, determination,
relation, or manifestation of something.

Change must be examined in relation to identity and continuity,
rather than being treated as automatically incompatible with them.


\subsection*{Continuity}

A form of persistence or connectedness through which identity,
determination, or manifestation may remain intelligible across
change.

Its precise ontological conditions require systematic examination.


\subsection*{Invariance}

A condition or relation of preservation across specified variations
or transformations.

Invariance should not be confused with absolute immobility or
the complete absence of change.


\subsection*{Persistence}

The continuation of something across a period or through a sequence
of conditions.

Persistence alone does not establish the complete nature of identity
or the conditions of viability.


\subsection*{Transformation}

A process through which the condition, structure, relation, or
manifestation of something is altered.

The relationship between transformation and identity requires
ontological analysis.


% =============================================================================
% PART VI — TERMINOLOGICAL AND METHODOLOGICAL DISTINCTIONS
% =============================================================================

\section*{VI. Terminological and Methodological Distinctions}


\subsection*{Ontological}

Pertaining to being, existence, identity, and the conditions or
structures of reality.


\subsection*{Epistemological}

Pertaining to knowledge, knowing subjects, the conditions of knowing,
and the limits or structure of knowledge.


\subsection*{Ontological Priority}

A claim concerning the foundational or constitutive status of one
ontological condition in relation to another.

No such priority should be inferred merely from the order in which
concepts are introduced or discussed.


\subsection*{Conceptual Distinction}

A distinction made between concepts for purposes of analysis.

A conceptual distinction does not necessarily imply that the concepts
refer to ontologically independent entities.


\subsection*{Ontological Identity and Existence}

The inseparable relation between identity and existence within
the present philosophical framework.

Neither is treated as ontologically prior to the other.


\subsection*{Open Philosophical Question}

A question whose definitive answer has not yet been established
within the current state of the philosophical system.

An open question must not be presented as a settled axiom,
definition, or thesis.


\subsection*{Provisional Definition}

A working definition adopted to maintain terminological consistency
while the relevant philosophical investigation remains incomplete.

A provisional definition may be revised as the system develops.


% =============================================================================
% PART VII — STATUS AND USE OF THIS GLOSSARY
% =============================================================================

\section*{VII. Status and Use of This Glossary}


The entries in this glossary have different degrees of conceptual
determination. Their inclusion does not imply that every term has
the same degree of philosophical finality.

The following editorial distinctions govern their interpretation:

\begin{description}

    \item[Established terminology:]
    Terms whose basic meaning is supported by the current
    foundational commitments of the project.

    \item[Structural terminology:]
    Terms used to identify components of the current architectural
    organization of the philosophical system.

    \item[Provisional terminology:]
    Terms whose working meaning is provided for consistency,
    while their complete ontological determination remains open.

    \item[Open questions:]
    Concepts or relations for which no definitive ontological
    position is adopted at this stage.

\end{description}


The glossary is therefore a controlled and revisable terminological
reference. Its entries must remain consistent with the arguments,
distinctions, and formal commitments established in the main text.

Where a detailed chapter develops a more precise definition,
the chapter's argument governs the interpretation of the term.
The glossary must subsequently be revised to reflect that development.


% -----------------------------------------------------------------------------
% End of File
% -----------------------------------------------------------------------------
